\section{A hyperbolic example}\label{sec:hyperbolic}

Fix any prime \(p\), and use the base \(G=G_p=\pi_1X\) constructed
above over \(\F_Q\), with
\begin{equation}\label{eq:hyp-parameters}
 \begin{gathered}
 g=100,\quad d=103,\quad c=2060^{-103},\quad
 r=100c^{-1},\quad n=r-1,\\
 Q=p^a>320000\cdot2060^{309},\quad m=Q^d,\quad
 N\ge cm/2,\quad |\cE|\le m/(2r).
 \end{gathered}
\end{equation}
In particular \(n\ge2\) and \(N\ge2\). We use the ceiling of
Proposition~\ref{prop:rank-budget} for every unital homomorphism
\(\F_p[G]\to D\), without requiring that homomorphism to be injective.
All geometry below uses only the incidence girth and the fixed triangle
shape, and is independent of \(p\).

\subsection{The permutation graph of groups}

Let \(U=\{1,\ldots,n\}^2\) and \(\Pi=\operatorname{Sym}(U)\).
For every \(i\in I\) let \(A_i\) be free on
\(\{a_{i,uv}:(u,v)\in U\}\), and put \(A=*_{i\in I}A_i\).
Take one central vertex with group \(A\) and, for every \(\pi\in\Pi\),
one vertex with a marked copy \(G^\pi\) of \(G\). For each pair
\((\pi,i)\) place one edge from \(G^\pi\) to \(A\), with edge group
\(F_i\), and edge maps
\begin{equation}\label{eq:hyp-edge-maps}
 e_{i,uv}\longmapsto e_{i,uv}^{\pi}\in G^\pi,
 \qquad e_{i,uv}\longmapsto a_{i,\pi(u,v)}\in A_i\le A.
\end{equation}
The notation \(\pi(u,v)\) denotes a single entry position, with no
requirement to preserve rows or columns. Both maps are injective by
Lemma~\ref{lem:free-core} and the displayed free bases.
Let \(\Gamma_p\) be the fundamental group of this graph of groups.
Parallel edges are intentional: there are \(1+(n^2)!\) vertices and
\(N(n^2)!\) edges.

Choose a line \(i_*\), and take all edges \((\pi,i_*)\) as a maximal
tree. The graph-of-groups presentation has one element \(t_{\pi,i}\)
per edge, with \(t_{\pi,i_*}=1\), and relations
\begin{equation}\label{eq:hyp-stable-letters}
 t_{\pi,i}e_{i,uv}^{\pi}t_{\pi,i}^{-1}=a_{i,\pi(u,v)}
 \qquad ((u,v)\in U).
\end{equation}
There is one common conjugator for the entire basis of an edge group.
Bass--Serre normal form embeds every vertex group and supplies the
action on its tree without inversions
\cite[Theorem~11 and its Corollary~3 in I.5.2, and Theorem~12 in I.5.3]{SerreTrees}. In normal-form
terms, choose coset representatives for each incident edge image:
an edge reversal can cancel only when its intervening vertex syllable
lies in that image. A nonidentity word within a vertex cannot reduce
to the empty path. This explains why graph cycles cause no extra
identifications within a vertex group.

A finite-order element fixes a vertex of the Bass--Serre tree. Indeed,
the finite hull of its finite orbit is invariant, and its center is a
vertex or an edge; an edge center also gives fixed endpoints because
there are no inversions. This is also~\cite[Proposition~19 in I.4.3]{SerreTrees}.
The vertex groups are torsion-free, so the element is trivial. Thus \(\Gamma_p\) is torsion-free.
It is finitely presented: take a finite presentation for each vertex,
\((N-1)(n^2)!\) stable letters, and the \(n^2\) relations
\eqref{eq:hyp-stable-letters} for each edge. This is finite and fixes the
construction from the deterministic base data.

\subsection{Permutation averaging over a division ring}

\begin{lemma}[Permutation lemma]\label{lem:permutation}
Let \(b_1,\ldots,b_{n^2}\) be pairwise distinct nonzero elements of
any division ring \(D\), with \(n\ge2\). If \(B^\pi\) is their
uniformly permuted \(n\times n\) array, then
\begin{equation}\label{eq:permutation-bound}
 \frac1{(n^2)!}\sum_\pi\rk_D B^\pi
 \ge(n-1)\left(1-\frac1{2n}\right)\ge\frac{n-1}{2}.
\end{equation}
\end{lemma}

\begin{proof}
First consider bordering an arbitrary \((k-1)\)-square matrix \(T\):
\[
 \begin{pmatrix}T&u\\v&z\end{pmatrix}.
\]
Invertible row and column operations put \(T\) in the form
\(\diag(I_s,0)\). Apply the same operations to \(u,v\), writing their
parts as \(u_1,u_0\) and \(v_1,v_0\). Eliminate the identity block.
The remaining matrix is
\[
 \begin{pmatrix}0&u_0\\v_0&z-v_1u_1\end{pmatrix}.
\]
If \(u_0\) or \(v_0\) is nonzero, its rank is positive for every
\(z\). Otherwise its rank is zero for at most the one value
\(z=v_1u_1\). Thus, with all other entries fixed, bordering fails to
increase rank for at most one corner value. The order of multiplication
is fixed by elimination; no commutative determinant is used.

Expose a uniform array by successively revealing the leading
\(k\times k\) square, for \(k=1,\ldots,n-1\). At each stage reveal its
new row and column before the corner. At that point \(k^2-1\) entries
are known, so the corner is uniform among exactly
\(n^2-k^2+1\) unused distinct values. Conditional on these exposed
entries, the probability of a rank increase is at least
\(1-1/(n^2-k^2+1)\). This also holds for \(k=1\), when the actual
increase is certain because the entries are nonzero. Bordering never
decreases rank. Each unused count is at least \(2n\), because
\(k\le n-1\). Adding the expected increases and ignoring the final
border gives
\[
 \mathbb E\rk_D B^\pi
 \ge\sum_{k=1}^{n-1}\left(1-\frac1{n^2-k^2+1}\right)
 \ge(n-1)\left(1-\frac1{2n}\right).
\]
This proves both bounds. There is no conditioning on previous pivot
success, and no independence between stages is needed.
\end{proof}

Suppose now that \(\phi:\F_p[\Gamma_p]\to D\) is injective.
For fixed \(i\), the elements
\[
 b_{i,uv}=\phi(a_{i,uv})-1\qquad ((u,v)\in U)
\]
are nonzero and pairwise distinct. The central vertex embeds, its
basis elements are distinct nonidentity group elements, and \(\phi\)
preserves the nonzero algebra elements \(a_{i,uv}-1\) and
\(a_{i,uv}-a_{i,u'v'}\), also in characteristic two.
Put \(B_i^\pi=(b_{i,\pi(u,v)})_{u,v=1}^n\).
The diagonal normalization after Lemma~\ref{lem:free-core}, followed
by subtracting row zero from the other rows and clearing row zero,
gives
\[
 Y_i\sim\begin{pmatrix}1&\one^{\mathsf T}\\
                 \one&(e_{i,uv})\end{pmatrix}
       \sim\diag\bigl(1,(e_{i,uv}-1)\bigr).
\]
By~\eqref{eq:hyp-stable-letters}, the copied difference matrix is
\((\phi(t_{\pi,i})^{-1}I)B_i^\pi(\phi(t_{\pi,i})I)\).
These scalar diagonal matrices are invertible even when \(D\) is
noncommutative. Hence
\begin{equation}\label{eq:hyp-rank-match}
 \rk_D\phi(Y_i^\pi)=1+\rk_D B_i^\pi.
\end{equation}
Restriction to each copy \(G^\pi\) satisfies
Proposition~\ref{prop:rank-budget}, with all the constants
in~\eqref{eq:hyp-parameters}:
\begin{equation}\label{eq:hyp-upper}
 \sum_i\rk_D\phi(Y_i^\pi)
 \le2(m-|\cE|)+r|\cE|
 \le\frac{5m}{2}<\frac{25m}{2}\le\frac{Nr}{4}.
\end{equation}
But the finite average, using the same \(\pi\) for every \(i\), obeys
\begin{equation}\label{eq:hyp-lower}
 \frac1{|\Pi|}\sum_{\pi\in\Pi}\sum_i\rk_D\phi(Y_i^\pi)
 \ge N\left(1+\frac{n-1}{2}\right)=\frac{Nr}{2}.
\end{equation}
Linearity of this average needs no independence between the line
arrays. Equations~\eqref{eq:hyp-upper}--\eqref{eq:hyp-lower} contradict
one another. Thus \(\F_p[\Gamma_p]\) embeds in no division ring.
The copies indexed by all \((n^2)!\) permutations are used only for
this averaging; the full symmetric group gives the simplest uniform
family of placements.

\subsection{Convex line lifts and joint malnormality}

\begin{lemma}\label{lem:line-geometry}
Each \(F_i\) is quasiconvex in \(G\), and the family is jointly
malnormal:
\begin{equation}\label{eq:joint-malnormal}
 F_i\cap gF_jg^{-1}\ne1\quad\Longrightarrow\quad i=j\text{ and }g\in F_i.
\end{equation}
\end{lemma}
\begin{proof}
The proof of Lemma~\ref{lem:free-core} gives the stronger fact that
the old link distance between distinct same-label \(x\)-germs at a
tip is greater than \(g\theta>2\pi\), or infinite. Every reduced
path in \(\Theta_i\) is therefore locally geodesic. Any path in its
universal-cover tree maps to a local geodesic in \(\widetilde X\),
hence to a geodesic~\cite[Proposition~II.1.4(2)]{BH}.
The tree embeds isometrically. Its image \(T\) is convex, since the
unique geodesic between any two of its points is that tree segment.
Its stabilizer is the corresponding conjugate of \(F_i\), acting
cocompactly. Every point of \(T\) lies at a uniformly bounded distance
from an \(F_i\)-orbit. Convexity and the orbit quasi-isometry
\cite[Proposition~I.8.19]{BH}, together with stability of
quasigeodesics~\cite[Theorem~III.H.1.7]{BH}, give quasiconvexity in a
Cayley graph of \(G\).

Distinct lifts of one \(\Theta_i\) are disjoint components of its
preimage; lifts with different labels have no common edge. If a
nonidentity \(h\) stabilized two such trees, its free action on each
tree would be a translation with an axis. There are no inversions,
since a deck transformation cannot fix an edge midpoint. Each tree
axis is a geodesic line in \(\widetilde X\) by convexity. The two
\(h\)-axes are asymptotic, so the flat strip theorem
\cite[Theorem~II.2.13]{BH} makes them bound a flat strip. A
\(\CAT(-1)\) space contains no flat strip of positive width, so they
coincide. This common line contains edges, contradicting the stated
intersection properties unless the two lifts are identical and have
the same label. Stabilizers of lifts, and the cosets indexing them,
now give~\eqref{eq:joint-malnormal}.
\end{proof}

\subsection{A sharper metric diffuseness lemma}

The systole estimate below is \(2\), so we need a bound sharper than
Lemma~\ref{lem:diffuse-metric}. For hyperbolic space the sharp form
is due to Bowditch: the fundamental group of a complete hyperbolic
manifold with injectivity radius greater than \(\log(1+\sqrt2)\),
so that every translation length exceeds
\(2\log(1+\sqrt2)=\arccosh3\), is diffuse
\cite[Theorem~5.3 and Lemma~5.4]{Bowditch}. The following lemma
gives the same threshold for free actions on \(\CAT(-1)\) spaces.
\begin{lemma}[Midpoint criterion]\label{lem:midpoint-diffuse}
Let a group act freely by isometries on a \(\CAT(-1)\) space. If every
nonidentity element has translation length at least \(\arccosh3\),
then the group is diffuse.
\end{lemma}
\begin{proof}
Fix \(\gamma\ne1\) and \(x\), and write
\(d=d(x,\gamma x)\), \(e=d(x,\gamma^2x)\). If \(z\) is the midpoint
of \([x,\gamma x]\), then \(\gamma z\) is the midpoint of
\([\gamma x,\gamma^2x]\). For the translation length \(L\), comparison
in this isosceles triangle gives
\[
 \cosh L\le\cosh d(z,\gamma z)
 \le1+\frac{\cosh e-1}{2(\cosh d+1)}.
\]
Indeed, the model midpoints have distance \(d/2\) from the common
vertex; the hyperbolic cosine law gives the displayed expression.
Thus \(\cosh L\ge3\) implies
\[
 \cosh e\ge4\cosh d+5>\cosh(d+\log4),\qquad e-d>\log4.
\]
For the strict comparison, use
\(\cosh(d+\log4)=(17\cosh d+15\sinh d)/8<4\cosh d\).
If both \(\gamma x\) and \(\gamma^{-1}x\) were no farther from some
\(o\) than \(x\), the four-point inequality~\eqref{eq:four-point},
from~\cite[Definition~4.1 and Theorem~5.1]{NicaSpakula}, would give
\[
 e^{(d(o,x)+e)/2}\le2e^{(d(o,x)+d)/2},
\]
contrary to \(e-d>\log4\). At least one neighboring translate is
therefore farther from \(o\). Two diameter endpoints of any finite
subset of a free orbit are extremal, using the other endpoint as
\(o\). This proves diffuseness without a semisimplicity hypothesis.
\end{proof}

\subsection{Separate cones and diffuse quotients}

Fix one label \(i\). Attach the geometric cone on \(\Theta_i\) to
\(X\), with a new apex \(U_i\), radial sides of length \(3\), and
the same right isosceles hyperbolic triangles along its \(x_{iuv}\)
edges. Denote this finite complex by \(Z_i\). Only this one line is
coned; the construction is performed separately for each \(i\).
Set
\[
 L_i=\langle\!\langle F_i\rangle\!\rangle_G,\qquad
 Q_i=G/L_i=\pi_1Z_i.
\]

\begin{lemma}\label{lem:cone-systole}
The universal cover \(\widetilde Z_i\) is proper \(\CAT(-1)\), every
essential loop in \(Z_i\) has length at least \(2\), and \(Q_i\) is
diffuse.
\end{lemma}
\begin{proof}
The link at the new apex is \(K_{r,r}\) with edges \(\pi/2\), so its
circuits have length at least \(2\pi\). An old tip link gains one
radial vertex, with \(\theta\)-edges to the same-\(i\) germs.
Any new simple circuit uses that vertex once, and its old part joins
two distinct same-label germs. Lemma~\ref{lem:free-core} gives old
length greater than \(g\theta\), so the new circuit has length greater
than \(2\theta+g\theta>2\pi\). This uses the full link separation,
not merely local geodesicity of line paths. The old apex links are
unchanged. At an edge interior the link consists of semicircles
between its two tangent directions, so all circuits have length at
least \(2\pi\); face interiors cause no difficulty. The finite-shape
link condition and Cartan--Hadamard
\cite[Theorems~I.7.50, II.5.2 and II.4.1(2), Lemma~II.5.6]{BH}
give the asserted universal cover. It is proper by local finiteness.

Suppose an essential loop has length less than \(2\). As in
Proposition~\ref{prop:systole}, the radius-\(3\) ball about any tip
lies in its contractible open star; the new triangles have the same
shape and do not change that argument. If the loop came within
\(2\) of a tip, shorter arcs of length less than \(1\) would put the
entire loop in that star. Thus all transition points on \(x\)-edges
have along-edge distances at least \(2\) from both tips, and signed
midpoint coordinate \(|s|\le a/2-2\le1\).

Use a finite string of cell segments and reduce its block--edge record
exactly as in Proposition~\ref{prop:systole}, now including the cone
as a block \(c\). A block with the tips removed is contractible and
lifts to \(\widetilde Z_i\). Returning through the same \(x\)-edge
can be replaced by the geodesic edge segment, without increasing
length; its endpoints and the replacement remain in the same
midpoint interval. Combining neighboring passages decreases their
number, so the process terminates. A disappearing record or a record
contained in one block contracts. Otherwise every passage joins
distinct \(x\)-edges and consecutive blocks are distinct.

For old blocks, the two estimates in
Proposition~\ref{prop:systole} still apply to these lifts: any passage
costs at least \(\eps>79/1000\), and a same-line passage costs at
least \(2h>3/2\). For a new-cone passage, we claim a cost greater
than \(1/2\). If its two edges share a tip, their angle there is
\(2\theta\): the new two-edge link path is shorter than any old
path, whose length exceeds \(g\theta\). Both along-edge distances
from that tip are at least \(2\). Hinge comparison gives
\begin{equation}\label{eq:cone-hinge}
 D\ge2\arcsinh(\sinh2\sin\theta)>\frac12.
\end{equation}
Here the minimum over two unequal radial distances also occurs at
\((2,2)\). To verify this, let \(c_1,c_2\) be the model rays at angle
\(2\theta\). The function \(f(s,t)=d(c_1(s),c_2(t))\) is jointly
convex by convexity of distance~\cite[Proposition~II.2.2]{BH}.
The cosine law gives both partial derivatives of \(\cosh f\) at
\((2,2)\) as \(\sinh2\cosh2(1-\cos2\theta)>0\).
The supporting-plane inequality for this differentiable convex
function therefore minimizes \(f\) on \([2,\infty)^2\) at \((2,2)\).
For the strict numerical bound, \(\sinh2>7/2\),
\(\sin\theta>1/12\), and
\(\sinh(1/4)<(1/4)/(1-1/16)=4/15<7/24\).

If the edges share no tip, their directions at the cone apex have
link distance at least \(\pi/2\), and their radial distances are at
least \(h\). Hinge comparison yields
\(\cosh D\ge\cosh^2h>\cosh(1/2)\), again giving the claimed cost.
For example, \(h>3/4\) gives \(\cosh^2h>25/16\), whereas the series
for \(\cosh(1/2)\) bounds it above by \(8/7\). All these are bounds
on endpoint distance in the universal cover and so apply to passages
crossing several cells as well.

Put \(\delta=79/1000\). If there are \(k\) passages, each has length
greater than \(\delta\), whence \(2k<4/\delta=4000/79<100=g\).
Project \((j,u,v)\) to its line label \(j\). The augmented block--line
graph is the old graph with a single leaf \(c\) attached only to
line \(i\); its girth is still greater than \(100\), and every visit
to \(c\) is a reversal. The image of the projected closed walk,
whose length is \(2k\), is a finite tree with an edge. A line vertex
cannot be a leaf of this image, since that would repeat consecutive
blocks. Thus at least two leaves are block vertices, and at most one
is \(c\). Two old leaf reversals cost at least \(4h>3\). One old
reversal and a cone reversal cost more than \(2h+1/2>2\).
Both contradict the loop length. This proves the systole bound.

Every nonidentity deck displacement projects to an essential loop
and is at least \(2>\arccosh3\). The deck action is free, so
Lemma~\ref{lem:midpoint-diffuse} proves \(Q_i\) diffuse.
\end{proof}

\subsection{Punctures and the free-product kernel}

We give the topological calculation in full. The connecting paths in
its statement are chosen as part of the conclusion; arbitrary choices
of paths are not asserted to yield a free-product isomorphism.

\begin{lemma}[Puncture lemma]\label{lem:punctures}
Let \(M\) be a proper, locally finite \(\CAT(0)\)
\(M_\kappa\)-polyhedral complex, with \(\kappa\le0\) and finitely
many shapes. Let \(P\) be a locally finite set of vertices whose links
are connected. Then \(M\setminus P\) is path-connected. There exist
connecting paths \(\gamma_v\) from a common basepoint to small link
spheres about \(v\in P\) for which the inclusion homomorphisms give
an isomorphism
\begin{equation}\label{eq:puncture-product}
 \mathop{*}_{v\in P}\pi_1(\operatorname{Lk}_M v)
 \ \xrightarrow{\ \cong\ }\ \pi_1(M\setminus P),
 \qquad [\alpha]\longmapsto[\gamma_v\alpha\gamma_v^{-1}].
\end{equation}
\end{lemma}
\begin{proof}
If \(P\) is empty, \(M\) is contractible and the assertion follows.
Otherwise choose \(p\in P\), put \(B_R=\overline B(p,R)\) and
\(Y_R=B_R\setminus P\), and start with a basepoint on a small sphere
about \(p\). Properness and local finiteness imply that the distance
set \(\{d(p,v):v\in P\}\) meets bounded intervals in finite sets,
with finite multiplicity. Small neighborhoods of vertices are balls
in the \(\kappa\)-cone on their links
\cite[Theorem~I.7.39]{BH}. Geodesics are unique and depend continuously
on their endpoints~\cite[Proposition~II.1.4(1)]{BH}.

\smallskip\noindent
\emph{Bands without punctures.}
If \([R,R']\) contains no puncture distance, radially move \(x\in
Y_{R'}\) along \([p,x]\) to radius \(\min\{d(p,x),R\}\).
This is a strong deformation retraction to \(Y_R\). Only radii in
\([R,R']\) move, so no puncture is met. For a sufficiently small
positive \(R\), \(Y_R\) is a punctured cone ball at \(p\); it retracts
to \(\operatorname{Lk}_M p\), which is connected.

\smallskip\noindent
\emph{Crossing a puncture level.}
Let \(v_1,\ldots,v_k\) be the punctures at distance \(d>0\).
Choose \(\delta>0\), with \(d-\delta>0\), so small that
\([d-3\delta,d+3\delta]\) contains no other puncture distance,
the \(2\delta\)-balls about these vertices are pairwise disjoint
cone neighborhoods, and they contain no other punctures. Write
\(W=Y_{d+\delta}\). Define the shadow
\[
 S_j=\{x\in B_{d+\delta}:v_j\in[p,x]\}.
\]
It is closed, since geodesics vary continuously, and it is contained
in \(\overline B(v_j,\delta)\), since
\(d(p,x)=d+d(v_j,x)\) on the shadow.
If \(\xi_p\) is the direction from \(v_j\) toward \(p\), cone
coordinates \((u,t)\) satisfy the shadow criterion
\begin{equation}\label{eq:shadow-criterion}
 v_j\in[p,x]\quad\Longleftrightarrow\quad
 \angle_{v_j}(p,x)=\pi
 \quad\Longleftrightarrow\quad
 d_{\operatorname{Lk}v_j}(\xi_p,u)\ge\pi.
\end{equation}
For the forward implication, the two subsegments form a geodesic.
For the converse, angle \(\pi\) makes their concatenation locally
geodesic and hence geodesic in a \(\CAT(0)\) space
\cite[Proposition~II.1.4(2)]{BH}. The angle is the link distance capped
at \(\pi\) in a cone neighborhood. This also follows from the angle
comparison formulation~\cite[Proposition~II.1.7(4)]{BH}.

Set \(O=W\setminus\bigcup_jS_j\). The radial homotopy toward radius
\(d-\delta\) stays in \(O\): if \(y\in[p,x]\) belongs to \(S_j\),
then \(x\) belongs to \(S_j\). It meets no earlier puncture, since
only radii at least \(d-\delta\) move, and it avoids the present ones
by removing their shadows. Thus \(O\) strongly deformation retracts
to \(Y_{d-\delta}\).

Put \(V_j=W\cap B(v_j,2\delta)\). These sets and \(O\) are open
relative to \(W\); the \(V_j\)'s are disjoint, and together they
cover \(W\). In cone coordinates about \(v_j\), every ray meets
\(V_j\) in an initial interval containing \((0,\delta]\): the bound
\(d(p,(u,t))\le d+t\) includes this interval, and convexity of
distance along the ray~\cite[Proposition~II.2.2]{BH} makes the entire
sublevel interval initial. Choose \(0<\eta<\delta\). Radial
interpolation of \(t\) to \(\eta\) stays in these intervals, giving
\(V_j\simeq\operatorname{Lk}_M v_j\).

Inside \(V_j\), no other shadow occurs, so
\eqref{eq:shadow-criterion} identifies \(O\cap V_j\) with the rays
having \(d_{\operatorname{Lk}v_j}(u,\xi_p)<\pi\), with those same
initial radial intervals. It retracts onto
\[
 B_{\operatorname{Lk}v_j}(\xi_p,\pi)\times\{\eta\}.
\]
Links are \(\CAT(1)\); this nonempty open \(\pi\)-ball is
contractible by contraction to its center along unique continuously
varying geodesics~\cite[Proposition~II.1.4(1),(4)]{BH}.
The overlap is therefore path-connected and simply connected.
Starting with connected \(Y_{d-\delta}\), the cover proves that \(W\)
is connected. Van Kampen, applied to the disjoint \(V_j\)'s one at a
time, gives
\[
 \pi_1(W)\cong\pi_1(Y_{d-\delta})*
       \left(\mathop{*}_{j=1}^k\pi_1(\operatorname{Lk}_M v_j)\right).
\]
The factor maps are the actual link inclusions with connecting paths
through \(O\); these paths may be chosen to extend all earlier ones.

\smallskip\noindent
\emph{Assembly.}
The band retractions and these free-product extensions exhaust
\(M\setminus P\). Every loop and every null-homotopy has compact
image and is contained in some \(Y_R\). Thus its fundamental group
is the direct limit of the displayed groups, and its path components
are already accounted for by the connected stages. The inclusions
are injective free-factor maps; the limit gives exactly
\eqref{eq:puncture-product} with the chosen paths.
\end{proof}

Apply Lemma~\ref{lem:punctures} to \(M=\widetilde Z_i\) and to the
lifts \(P\) of the new apex. Their links are copies of
\(\Theta_i=K_{r,r}\). The inverse image of \(X\) is connected,
because \(G\to Q_i\) is onto, and is the cover \(X_{L_i}\)
corresponding to \(L_i\). Each lifted cone embeds (its domain is a
cone and its base vertices and edges are distinct), and a punctured
cone retracts radially to its base. These retractions agree on
\(X_{L_i}\); local finiteness makes their union a deformation
retraction of \(M\setminus P\) onto \(X_{L_i}\).
The apex lifts are indexed by \(Q_i\), and the link factors become
actual conjugates of the line subgroup. Consequently, for suitable
representatives and connecting paths,
\begin{equation}\label{eq:line-kernel-product}
 L_i=\mathop{*}_{q\in Q_i}g_qF_i g_q^{-1}.
\end{equation}
All factors are free, so \(L_i\) is a free group, possibly of infinite
rank. In particular one conjugate of \(F_i\) is a free factor of \(L_i\).
Conjugating the whole decomposition by its \(g_q^{-1}\) preserves
\(L_i\), since \(L_i\lhd G\), and makes \(F_i\) itself a literal
free factor. This is the required kernel statement; contractibility
alone would not supply it.

\subsection{Concentration and the two combination steps}

For a group \(T\) acting on a set, a point \(x\) of a finite set
\(S\) is extremal if \(\gamma x,\gamma^{-1}x\in S\) forces
\(\gamma x=x\). The action is diffuse when every finite subset
with at least two points has at least two extremal points.
A subgroup \(H\le T\), for \(T\) torsion-free, is
\emph{concentrated} when the left-coset action on \(T/H\) is diffuse.
If \(L\lhd T\), then \(L\) is concentrated in \(T\) exactly when
\(T/L\) is diffuse: the action factors through the regular action of
\(T/L\), and the extremal-point condition is identical. This proves
the normal-subgroup remark in~\cite[p.~4 of the author version]{Bowditch}.

Every free group, including one of infinite rank, is diffuse by its
free action on a tree~\cite[Theorem~1.2(3)]{Bowditch}. Every free factor
\(H\) of a free group \(H*H'\) is concentrated: apply
\cite[Theorem~1.3(2)]{Bowditch} to this splitting with trivial edge
group. Its edge group is diffuse and is concentrated at each end
because both vertex groups are diffuse; \(H\) is concentrated in
itself, so the theorem's additional conclusion concentrates it in
\(H*H'\). This applies also when \(H'\) has infinite rank.

By~\eqref{eq:line-kernel-product}, \(F_i\) is concentrated in \(L_i\).
By Lemma~\ref{lem:cone-systole}, \(Q_i\) is diffuse, so \(L_i\) is
concentrated in \(G\). Transitivity of concentration
\cite[Theorem~1.3(1)]{Bowditch} now proves that each \(F_i\) is
concentrated in \(G\). This is a separate conclusion of the cone
and puncture arguments, not a consequence of quasiconvexity and
malnormality alone.

Each edge group of the permutation graph is intrinsically free and
diffuse, concentrated in \(G^\pi\) by the preceding argument, and
concentrated in \(A\) because \(A_i\) is a free factor. We already
proved \(\Gamma_p\) torsion-free. Bowditch's graph-of-groups theorem
\cite[Theorem~1.3(2)]{Bowditch} therefore gives diffuseness of
\(\Gamma_p\). Its hypotheses are per incident edge: the incident
edge set at a vertex is a disjoint union of diffuse coset sets.
Repeated images \(A_i\) and cycles in the underlying graph are
allowed, as is explicit in that theorem and its proof.

For hyperbolicity, \(A\) and every \(G^\pi\) are word-hyperbolic.
The edge images are quasiconvex, by Lemma~\ref{lem:line-geometry} on
the \(G\) side and convexity of a free factor in the basis tree on
the \(A\) side. They are proper: \(N\ge2\) makes \(A_i<A\), and
\(F_i=G\) would contradict joint malnormality with a nontrivial
\(F_j\), \(j\ne i\). All vertex and edge groups are non-elementary,
since they contain nonabelian free groups.

At a lifted \(G^\pi\)-vertex, edges are indexed by pairs
\((i,gF_i)\). Two distinct such edges have trivially intersecting
stabilizers by~\eqref{eq:joint-malnormal}, including equal labels
with different cosets. Every three-edge segment in the bipartite
Bass--Serre tree has an internal \(G^\pi\)-vertex. Therefore
\begin{equation}\label{eq:three-acylindrical}
 \operatorname{Fix}_{\Gamma_p}(J)=1
 \quad\text{for every segment \(J\) of three edges}.
\end{equation}
At a central vertex, two distinct edges with the same \(i\) and the
same coset can have the common stabilizer \(aA_i a^{-1}\).
Thus two-edge stabilizers need not be trivial; equation
\eqref{eq:three-acylindrical} is the acylindricity needed here.

Apply Dahmani~\cite[Theorem~0.1(1), p.~935]{Dahmani} to this finite
acylindrical graph, viewing the hyperbolic vertex groups as relatively
hyperbolic with empty peripheral families. Proper quasiconvex edge
images are fully quasiconvex by~\cite[Proposition~1.11, p.~942]{Dahmani}.
The theorem gives relative hyperbolicity with peripheral family formed
from the vertex peripherals, hence empty. This is ordinary word
hyperbolicity; equivalently the resulting convergence action is
uniform, as explained in~\cite[p.~934]{Dahmani}, and the
uniform-convergence characterization applies
\cite[Theorem~0.1]{BowditchTop}. For related acylindrical combination results, see also
Kapovich~\cite{Kapovich}.

\begin{proof}[Completion of the proof of Theorem~\ref{thm:main}]
The permutation graph gives a finitely presented, torsion-free,
word-hyperbolic diffuse group \(\Gamma_p\). Diffuseness implies
unique product~\cite[Proposition~1.1]{Bowditch}; the support-coefficient
argument in Section~\ref{sec:diffuse} makes \(K[\Gamma_p]\) a domain
for every field \(K\). Equations~\eqref{eq:hyp-upper}--\eqref{eq:hyp-lower}
exclude an embedding of \(\F_p[\Gamma_p]\) into any division ring.
If \(\operatorname{char}K=p\), an embedding of \(K[\Gamma_p]\) would
restrict to one of its prime-field subalgebra \(\F_p[\Gamma_p]\).
The construction began with an arbitrary prime \(p\), so this proves
all the stated quantifiers.
\end{proof}
